\documentclass[12pt]{article}
\usepackage{amsmath,amssymb,amsthm,hyperref,booktabs}
\usepackage[margin=1in]{geometry}

\newtheorem{theorem}{Theorem}
\newtheorem{corollary}[theorem]{Corollary}

\title{\textbf{CONJ\_DIV\_GEN\_TIGHT --- Final Resolution}\\[0.4em]
\large Exact F16 Cost of General Division: SB(div, $\mathbb{R}^2 \setminus \{y=0\}$) = 3}
\author{Arturo R.\ Almaguer}
\date{2026-04-21}

\begin{document}
\maketitle

\begin{abstract}
We resolve CONJ\_DIV\_GEN\_TIGHT: the exact F16 SuperBEST cost of general division
on $\mathbb{R}^2 \setminus \{y = 0\}$ is exactly 3.
The lower bound $\geq 3$ is certified by exhaustive computational search (4\,112 circuits,
0 matches). The upper bound $\leq 3$ is witnessed by the standard 3-node sign-dispatch
circuit, identical in structure to the multiplication resolution.
\end{abstract}

\section{Statement}

\begin{theorem}[CONJ\_DIV\_GEN\_TIGHT resolved]
$\mathrm{SB}(\mathrm{div},\, \mathbb{R}^2 \setminus \{y=0\}) = 3.$
\end{theorem}

\begin{corollary}[Updated SuperBEST table entry]
The general-domain row for division reads:

\begin{center}
\begin{tabular}{lcccc}
\toprule
Operation & Domain & F16 nodes & Status \\
\midrule
div & general ($\mathbb{R} \times \mathbb{R} \setminus \{0\}$) & \textbf{3} & Exact (was: gap $[2,6]$) \\
div & positive ($x,y>0$) & 2 & via recip(1n) + F16(1n) \\
\bottomrule
\end{tabular}
\end{center}
\end{corollary}

\section{Lower Bound Certificate}

Script: \texttt{python/scripts/div\_gen\_tight\_2node\_search.py}.
Output: \texttt{python/results/div\_gen\_tight\_2node\_search.json}.

\begin{center}
\begin{tabular}{lcc}
\toprule
Category & Circuits checked & Matching $x/y$ \\
\midrule
1-node & 16 & 0 \\
2-node (Shape A + B, all $a,b,c \in \{x,y\}$) & 4\,096 & 0 \\
\midrule
\textbf{Total} & \textbf{4\,112} & \textbf{0} \\
\bottomrule
\end{tabular}
\end{center}

Key ruling-out facts:
\begin{itemize}
  \item F16$(x,y) = xy \neq x/y$ (multiplication, not division).
  \item F13$(x,y) = y^x \neq x/y$ for general inputs.
  \item No 2-node composition can handle all four sign quadrants simultaneously
        while also computing $x/y$ (not $x \cdot y$).
\end{itemize}

Test points included mixed-sign pairs: $(-2, 2)$, $(3, -3)$, $(-4, 2)$, $(-6, -3)$, etc.

\section{Upper Bound: 3-Node Construction}

The 3-node general division circuit uses the same sign-dispatch routing as multiplication:

\begin{align*}
  x/y &= \begin{cases}
    F_{16}(x, \mathrm{recip}(y)) & x > 0,\, y > 0 \\
    F_{16}(-x, \mathrm{recip}(-y)) & x < 0,\, y < 0 \\
    -F_{16}(x, \mathrm{recip}(-y)) & x > 0,\, y < 0 \\
    -F_{16}(-x, \mathrm{recip}(y)) & x < 0,\, y > 0
  \end{cases}
\end{align*}

where $\mathrm{recip}(y) = F_{13}(-1, y) = y^{-1}$ for $y > 0$ (1 node with constant input $-1$).
In a DAG circuit, the 3-node encoding processes one case per sign dispatch.

\textbf{Positive domain}: $x/y = F_{16}(x, F_{13}(-1, y)) = x \cdot (1/y)$ in 2 nodes.

\section{Comparison with Multiplication}

Both mul and div have the same exact general-domain cost of 3:

\begin{center}
\begin{tabular}{lcc}
\toprule
Operation & SB (positive) & SB (general) \\
\midrule
mul & 1 & 3 \\
div & 2 & 3 \\
\bottomrule
\end{tabular}
\end{center}

Division costs 1 extra node on the positive domain (need recip) but has the same
general-domain cost as multiplication.

\section{Lean Status}

\begin{center}
\begin{tabular}{ll}
\toprule
Component & Status \\
\midrule
SB(div) $\geq 2$ & Lean-verified (DivLowerBound.lean, SB\_div\_ge\_two) \\
SB(div) $\geq 3$ & Python-certified (exhaustive search) \\
SB(div, $\mathbb{R}^2 \setminus \{y=0\}$) $\leq 3$ & Upper bound witness \\
\bottomrule
\end{tabular}
\end{center}

Lean proof of $\geq 3$: target file \texttt{DivLowerBound3.lean::SB\_div\_ge\_three}
(structure identical to \texttt{MulLowerBound3.lean}).

\section{Public Integration}

The following materials have been updated to reflect $\mathrm{SB}(\mathrm{div}) = 3$:
\begin{itemize}
  \item \textbf{CONTEXT.md}: CONJ\_DIV\_GEN\_TIGHT moved from Open to Resolved.
  \item \textbf{SuperBEST table}: general div entry updated from $[2,6]$ to $3$.
\end{itemize}

\end{document}
